\section{Конструкторская часть}
\label{sec:design}

Конструкторская часть фиксирует архитектуру системы: какие модули в неё входят, как они
связаны и какие данные передают друг другу. Отправных точек две~--- выбранное
представление выражений (раскрученная двусторонняя очередь) и грамматика базисного
подмножества; от них исходит вся дальнейшая декомпозиция системы.

Начинается часть с требований~--- функциональных и нефункциональных,~--- а дальше
требования разворачиваются в цепочку этапов обработки, внутренние представления данных,
интерфейс контейнера выражений и, наконец, декомпозицию на модули с обработкой ошибок и
составом встроенной библиотеки. Связь требований с конкретными модулями держится в поле
зрения на всём протяжении раздела.

\subsection{Функциональные требования}

Функциональные требования следуют непосредственно из постановки задачи: система должна
принимать и исполнять программы на базисном подмножестве Рефала-5, сообщая об ошибках
входного текста. К системе предъявляются требования:
\begin{itemize}
  \item приём исходного текста на подмножестве Рефала-5;
  \item выдача сообщений об ошибках с привязкой к позиции во входном тексте;
  \item выполнение программы в модели \enq{выбор первого применимого правила} для функций;
  \item сокращённая библиотека встроенных функций (только необходимый минимум по постановке).
\end{itemize}

Из этих требований и складывается состав модулей. Приём текста и диагностика
ошибок ложатся на лексический и синтаксический анализаторы, хранящие координаты
токенов; модель «выбора первого применимого правила» воплощается в семантическом анализе
и нормализаторе; сокращённая библиотека встроенных функций задаёт границу того, какие
программы система обязана исполнять. Нефункциональные требования (переносимость,
воспроизводимость сборки) рассматриваются отдельно в п.~\ref{subsec:modules} и далее.

\subsection{Архитектура и потоки данных}

Система включает цепочку этапов: подготовка входа (при необходимости), лексический анализ, синтаксический анализ,
семантический анализ и выполнение правил. Представление объектных выражений едино для этапов выполнения и используется
при сопоставлении и сборке результата.

На практике удобно разрешить внешний \enq{рассахариватель}, который переводит исходный текст полного Рефала-5 к синтаксису
базисного подмножества (условия, блоки и др.), чтобы не расширять грамматику минимальной системы. Такой шаг не меняет
семантику рассматриваемых программ при условии корректности рассахаривания, но выносит сложные конструкции из ядра
транслятора. В реализации в качестве рассахаривателя используется внешний инструмент
\prog{r5fw-desugar} из проекта \prog{refal-5-framework}~\cite{refal5_framework_github};
сам интерпретатор его не запускает, но принимает на вход результат его работы (см.~п.~\ref{subsec:usage}).

\begin{figure}[!htb]
\centering
\begin{tikzpicture}[
    node distance=0.55cm and 0.55cm,
    box/.style={rectangle, draw, rounded corners, minimum width=2.35cm, minimum height=0.78cm, align=center, font=\footnotesize},
    data/.style={rectangle, draw, dashed, minimum width=2.15cm, minimum height=0.62cm, align=center, font=\scriptsize},
    arrow/.style={-{Stealth[length=2.2mm]}, thick}
]
    \node[box] (prep) {Подготовка\\входа};
    \node[box, right=of prep] (lexer) {Лексический\\анализ};
    \node[box, right=of lexer] (parser) {Синтаксический\\анализ};
    \node[box, right=of parser] (sem) {Семантический\\анализ};
    \node[box, right=of sem] (exec) {Исполнение\\правил};

    \node[data, below=0.55cm of prep] (srcfull) {Исходный текст};
    \node[data, below=0.55cm of lexer] (srcbasis) {Базисный текст};
    \node[data, below=0.55cm of parser] (toks) {Токены};
    \node[data, below=0.55cm of sem] (ast) {Дерево программы};
    \node[data, below=0.55cm of exec] (expr) {Выражения};

    \draw[arrow] (srcfull) -- (prep);
    \draw[arrow] (prep) -- (srcbasis);
    \draw[arrow] (srcbasis) -- (lexer);
    \draw[arrow] (lexer) -- (toks);
    \draw[arrow] (toks) -- (parser);
    \draw[arrow] (parser) -- (ast);
    \draw[arrow] (ast) -- (sem);
    \draw[arrow] (sem) -- (exec);
    \draw[arrow] (exec) -- (expr);
\end{tikzpicture}
\caption{Этапы обработки программы и основные данные}
\label{fig:arch}
\end{figure}

Целевой средой реализации выбрана платформа~JVM, конкретно --- Java~25~LTS~\cite{java_se25_spec}.
Это естественно сочетается с представлением программы и выражений как ссылочных структур и со
сборкой мусора, но накладные расходы GC должны учитываться при выборе гранулярности
объектов (термы, буферы раскрученной двусторонней очереди, узлы дерева программы).

\subsection{Внутренние представления данных}

Каждому этапу обработки соответствует своя структура данных, а граница между этапами~---
это смена представления. В проектировании фиксируются следующие структуры:
\begin{itemize}
  \item токены с координатами (строка/столбец) для точной диагностики;
  \item дерево программы как набор функций и правил сопоставления;
  \item объектные выражения как двусторонняя очередь термов.
\end{itemize}

Выбор каждой структуры продиктован требованиями. Координаты в токенах нужны, чтобы
сообщения об ошибках указывали точное место во входном тексте. Дерево программы хранит
функции и их предложения в форме, удобной для сопоставления образца с аргументом без
повторного разбора исходного текста. Объектные выражения вынесены в отдельный контейнер по другой причине: от стоимости
операций над ними (конкатенация, добавление и удаление на концах) зависит
производительность всего интерпретатора, поэтому им отведён отдельный интерфейс, описанный
далее.

\subsection{Интерфейс контейнера выражений}

Для независимости модулей выделяется интерфейс контейнера выражения, содержащий операции:
проверку пустоты, отделение терма с начала/конца, добавление/удаление на концах, итерацию и конкатенацию.
Конкретная реализация (мелкий и глубокий уровни раскрученной двусторонней очереди) скрыта за этим
интерфейсом и может заменяться без изменения логики сопоставления. В Java-реализации это
запечатанный интерфейс \prog{Expr<T>} с двумя реализациями~--- \prog{ShallowDeque<T>} и
\prog{DeepDeque<T>}; такая форма позволяет компилятору проверять полноту разбора случаев в
switch-выражениях и используется в реализации операции \prog{concat} (см.~раздел~\ref{sec:tech}).

\subsection{Декомпозиция на модули}
\label{subsec:modules}

С практической точки зрения систему удобно декомпозировать на набор модулей со строго
определёнными ответственностями. Для удобства соотнесения проекта и реализации модули
названы по соответствующим Java-пакетам:
\begin{itemize}
  \item \prog{refal.lexer} --- лексический анализатор: преобразование потока символов в поток
        токенов с координатами (\prog{Lexer}, \prog{Token}, \prog{TokenType}, \prog{LexException});
  \item \prog{refal.parser} --- синтаксический анализатор: построение дерева программы
        (\prog{Parser}, \prog{ParseException});
  \item \prog{refal.ast} --- определения узлов дерева программы (\prog{Term} как
        запечатанный интерфейс, наследники \prog{Sym}/\prog{Num}/\prog{Str}/\prog{Var}/\prog{Br}/\prog{Call},
        а также \prog{Rule}, \prog{Func} и \prog{Program});
  \item \prog{refal.semantic} --- семантический анализатор: проверка определённости функций,
        согласованности типов переменных и т.\,п. (\prog{SemanticChecker}, \prog{SemanticException},
        \prog{SemanticInfo});
  \item \prog{refal.deque} --- основная реализация контейнера выражения (\prog{Expr},
        \prog{ShallowDeque}, \prog{DeepDeque}, вспомогательный класс \prog{ExprOps}
        с реализацией конкатенации);
  \item \prog{refal.runtime} --- слой исполнения: сопоставление образца с выражением,
        подстановка результата, цикл нормализации, библиотека встроенных функций
        (\prog{Runtime}, \prog{Matcher}, \prog{Substitutor}, \prog{Bindings}, \prog{Builtins},
        \prog{BuiltinFn}, \prog{RefalException});
  \item \prog{refal.Main} --- интерфейс командной строки.
\end{itemize}

Семантический анализ выделен в отдельный этап, так как иначе диагностика ошибок класса \enq{вызов неопределённой
функции} либо \enq{связь \(s\)- и \(e\)-переменных под одним именем} приходилась бы на
стадию интерпретации и затрудняла локализацию причины ошибки.

Связанность при такой декомпозиции остаётся низкой: уточнение синтаксиса входного языка не задевает слой контейнера выражений,
а смена представления выражений~--- логику синтаксического анализа.

\subsection{Модель выполнения и точки расширения}

Модель выполнения фиксируется как \enq{вызов функции \(\rightarrow\) просмотр предложений сверху вниз \(\rightarrow\) первое успешное сопоставление}.
Отдельные точки расширения в проектировании:
\begin{itemize}
  \item добавление новых классов термов (при необходимости расширить подмножество);
  \item расширение механизма сопоставления (например, введение дополнительных ограничений или условий);
  \item расширение библиотеки встроенных функций без изменения ядра сопоставления;
  \item возможное добавление оптимизаций (индексация предложений по первым термам, мемоизация вызовов).
\end{itemize}

Сами оптимизации в задачи этого этапа не входят, но важно, что выбранный интерфейс
контейнера выражений их не блокирует: добавить индексацию предложений или мемоизацию можно
будет без переделки ядра.

\subsection{Нефункциональные требования}

Помимо функциональности, для системы важны следующие свойства:
\begin{itemize}
  \item предсказуемость потребления памяти: представление выражений должно использовать структурное разделение и избегать полного копирования;
  \item диагностика: сообщения об ошибках должны указывать координаты во входном тексте и причину ошибки;
  \item расширяемость: состав встроенных функций допускается сокращать/расширять без изменения ядра сопоставления;
  \item переносимость: выбор языка реализации со сборкой мусора и поддержкой
        обобщённых типов; в качестве конкретной платформы выбрана Java~25~LTS~\cite{java_se25_spec},
        что обеспечивает запуск на любой ОС с установленной JVM;
  \item воспроизводимость сборки: сборка и тестирование автоматизированы средствами
        Apache Maven~\cite{maven_guide}; зависимости (JUnit~5~\cite{junit5_guide},
        AssertJ~\cite{assertj_docs}, Checkstyle~\cite{checkstyle_docs},
        JaCoCo~\cite{jacoco_docs}) фиксируются версиями в \prog{pom.xml}.
\end{itemize}

\subsection{Обработка ошибок и форматы сообщений}

Ошибки целесообразно разделять на классы: лексические (неожиданный символ, незакрытый строковый литерал), синтаксические (неожиданная конструкция,
нарушение структуры определения функции), семантические (неопределённая функция, несовместимые типы переменных, вызов в образце,
отсутствие определения для \prog{\$ENTRY}) и ошибки времени выполнения (отсутствие применимого правила, превышение лимита шагов).
Для повышения качества диагностики в проектировании закладывается единый формат сообщения: координата (строка, столбец), тип ошибки
и краткое описание. Для каждой категории определён собственный класс исключения
(\prog{LexException}, \prog{ParseException}, \prog{SemanticException}, \prog{RefalException}),
сообщение которого начинается с координат вида \prog{line:col:}.

\subsection{Выбор и сокращение встроенной библиотеки}

Сокращённая библиотека встроенных функций проектируется как отдельный слой среды выполнения: каждая встроенная функция имеет строго определённый
интерфейс, а таблица примитивов отделена от механизма сопоставления. В Java-реализации интерфейс
\prog{BuiltinFn} имеет единственный метод
\prog{apply(Runtime rt, List<Term> args)}, возвращающий список термов (раскрытие реализовано
интерпретатором). Регистрация выполняется в \prog{Builtins.defaultBuiltins()}:
изменение или сокращение состава встроенных функций сводится к редактированию этой таблицы и
не затрагивает ядро сопоставления и нормализации.

\subsection{Представление термов и вложенность}

В выбранном подмножестве термы образуют минимально необходимый набор для выражения большинства \enq{структурных} преобразований:
атомарные термы (символы-слова, числа, строковые литералы) и скобочные термы. Скобочный терм содержит вложенное выражение, что фактически вводит
деревоподобную структуру данных поверх линейной последовательности. С точки зрения проектирования это означает, что контейнер выражения
должен хранить не только атомы, но и ссылки на вложенные выражения, причём вложенность может быть произвольной.

Здесь важно различать две роли контейнера:
\begin{itemize}
  \item на верхнем уровне контейнер обслуживает операции над последовательностью термов (концы, конкатенация, разбиение);
  \item на уровне скобок контейнер используется рекурсивно, так как содержимое скобок является самостоятельным выражением.
\end{itemize}
Одна и та же структура работает на всех уровнях вложенности одинаково~--- это упрощает
архитектуру и делает поведение системы предсказуемее.

\subsection{Окружение сопоставления и представление подстановок}

Результатом сопоставления образца и выражения является окружение (подстановка), связывающая переменные образца со значениями.
В проектировании выделяются три класса связываний:
\begin{itemize}
  \item \(s\)-переменная связывается с атомарным термом (число, символ-слово или строковый литерал);
  \item \(t\)-переменная связывается с термом (включая скобочный);
  \item \(e\)-переменная связывается с выражением (цепочкой термов).
\end{itemize}

С практической точки зрения удобно хранить подстановку в виде отображения \enq{имя переменной \(\rightarrow\) значение}, где значение
для \(e\)-переменных является контейнером выражения. Это позволяет строить результат без распаковки цепочек в списки/массивы:
результат собирается из фрагментов через операции конкатенации и добавления на концах, опираясь на интерфейс контейнера.
В Java-реализации эту роль выполняет класс \prog{Bindings} (см.~приложение~\ref{app:listings}).

\subsection{Требования к механизму сопоставления}

Чтобы разделить \enq{логику сопоставления} и \enq{операции над данными}, проектирование фиксирует требования к модулю сопоставления:
\begin{itemize}
  \item вход: образец (последовательность термов с переменными) и выражение-аргумент;
  \item выход: либо \enq{неуспех}, либо подстановка, корректная относительно ограничений типов переменных;
  \item свойство: если сопоставление успешно, то подстановка совместима со всеми повторными вхождениями переменных (одно имя --- одно значение).
\end{itemize}

Такая специфика позволяет в дальнейшем заменять стратегию перебора границ \(e\)-переменных (например, на более эффективную),
не меняя остальную архитектуру. При этом контейнер выражения должен обеспечивать операции, которые нужны сопоставлению: проверку равенства
фрагментов, получение префикса/суффикса и сборку новых выражений с разделением структуры.

\subsection{Сценарии использования системы}
\label{subsec:usage}

Проектирование исходит из основного сценария --- пакетный запуск: чтение программы из файла, запуск входной функции
с заданным выражением и вывод результата. В реализации этому соответствует команда

\begin{lstlisting}[language=bash,caption={Базовая схема запуска}]
java -jar target/refal5q.jar <file.ref>[+lib.ref...] [--entry Name] [--step-limit N] [--expr "args..."]
\end{lstlisting}

Для программ, использующих расширенный синтаксис Рефала-5 (условия \prog{,~e.Cond}, блоки
\prog{\{ : \}}), предусмотрен сценарий с предварительным запуском внешнего рассахаривателя
\prog{r5fw-desugar}~\cite{refal5_framework_github}: исходный файл сначала проходит через
рассахариватель, после чего полученный файл подаётся на вход интерпретатору. Сам
интерпретатор рассахариватель не запускает; такой подход выносит сложные синтаксические
преобразования за пределы ядра системы.

Интерактивный режим (последовательные запуски с разными входными выражениями без повторной загрузки
программы) в рамках данного проектирования не фиксируется; его можно рассматривать как возможную доработку или расширение --- например, для отладки
и исследования поведения правил.
